\section{Falsifiability without overclaiming}

The recovery replaces spectacular but weakly derived forecasts with a sequence
of gate-level falsifiers.

\subsection{Near-term theoretical falsifiers}

\begin{enumerate}
\item If no stable finite-density branch can yield the required phonon power
and sign, UVIR-001 rejects the proposed condensate origin of the local law.
\item If a normalized matter vertex cannot satisfy galactic and local bounds
within one action, MAT-001/SCR-001 reject the universal coupling architecture.
\item If no relativistic completion produces acceptable PPN and lensing
potentials, the local acceleration law cannot serve as a gravitational model.
\item If the periodic disk solution preserves the empirical tension of the
algebraic $C_{\rm obs}=2/3$ implementation, disk geometry cannot rescue that matching
hypothesis.
\item If derived reservoir, shape and wake sectors have no stable driven
anisotropic solution, CBR-002 rejects persistent topological Hubble anisotropy
within that architecture.
\end{enumerate}

\subsection{Observational claims held downstream}

NANOGrav spectral features require a derived eigenmode spectrum and detector
response.  Cluster offsets require causal wake dynamics, a collision
calculation and a lensing metric.  Stellar-population claims require an
independent star-formation model.  CMB and growth require a coherent background
and perturbation theory.  None is presented here as a clean framework-level
falsifier.

\subsection{Reproducibility rule}

Every future numerical claim must identify its source equations, code entry
point, data version, parameter inventory, validation controls and a
machine-readable summary.  A reported statistic must be explicitly computed
by the cited pipeline.  Negative gate results, including the free-field
$13/12$ result, remain part of the model record rather than being hidden by a
new ansatz.
